\section{First-order design}\label{sec:first-order}

This mirrors the derivation of Example 7.6 of the reader. The first-order
armature-to-speed relation is
\begin{equation}
  J\dot{\omega}=-b\omega+K_m u .
  \label{eq:motor-differential-equation}
\end{equation}
Its continuous-time transfer function is
\begin{equation}
  G(s)=\frac{\omega(s)}{u(s)}
      =\frac{K_m/J}{s+b/J}.
  \label{eq:motor-continuous-transfer}
\end{equation}
With zero-order hold sampling and \(T=\SI{0.08}{s}\), the discrete transfer
function is written as
\begin{equation}
  G(z)=\frac{b_0}{z+a_0},
  \label{eq:motor-discrete-transfer}
\end{equation}
where
\begin{equation}
  a_0=-\exp\!\left(-\frac{bT}{J}\right),
  \qquad
  b_0=\frac{K_m}{b}
      \left(1-\exp\!\left(-\frac{bT}{J}\right)\right).
  \label{eq:motor-discrete-parameters}
\end{equation}
The identified parameter vector is therefore
\begin{equation}
  x_e=\begin{bmatrix}a_0 & b_0\end{bmatrix}^{T}.
  \label{eq:identified-parameter-vector}
\end{equation}

The discrete plant model in \eqref{eq:motor-discrete-transfer} gives the
one-step regression
\begin{equation}
  y(k)=C(k)x_e(k),
  \qquad
  C(k)=\begin{bmatrix}-y(k-1) & u(k-1)\end{bmatrix}.
  \label{eq:rls-regression}
\end{equation}
The recursive least-squares update with forgetting factor \(\alpha\) is
\begin{align}
  L(k) &=
  \frac{P(k-1)C^{T}(k)}
       {C(k)P(k-1)C^{T}(k)+\alpha},
  \label{eq:rls-gain}\\
  x_e(k) &= x_e(k-1)+L(k)\left(y(k)-C(k)x_e(k-1)\right),
  \label{eq:rls-parameter-update}\\
  P(k) &= \frac{\left(I-L(k)C(k)\right)P(k-1)}{\alpha}.
  \label{eq:rls-covariance-update}
\end{align}
The initial values are
\begin{equation}
  P(0)=100I,
  \qquad
  x_e(0)=\begin{bmatrix}-0.5&1.0\end{bmatrix}^{T}.
  \label{eq:rls-initialisation}
\end{equation}
The covariance update is suspended without excitation, which
prevents covariance windup. The controller-design step uses a guard
\(|b_0|>10^{-3}\), avoiding division by a nearly singular plant gain.

The desired closed-loop polynomial is chosen as
\begin{equation}
  q(z)=z^2-0.5z+0.06=(z-0.3)(z-0.2).
  \label{eq:desired-polynomial}
\end{equation}
For \(T=\SI{0.08}{s}\), these discrete poles correspond to continuous poles
close to \(-15\,\si{s^{-1}}\) and \(-20\,\si{s^{-1}}\). The target settling
time is therefore about \(\SI{250}{ms}\). Writing \(q(z)=z^2+q_1z+q_0\), the
controller parameters are computed from the current plant estimates by
\begin{equation}
  d_0=\frac{q_0+a_0}{b_0},
  \qquad
  d_1=\frac{q_1+1-a_0}{b_0}.
  \label{eq:pole-placement-parameters}
\end{equation}
The implemented controller is the incremental PI law
\begin{equation}
  u(k)=u(k-1)+d_1e(k)+d_0e(k-1),
  \label{eq:incremental-pi}
\end{equation}
where \(e(k)=r(k)-y(k)\). The saturated command, not the unsaturated internal
value, is fed back as \(u(k-1)\) in \eqref{eq:incremental-pi}; this is the
anti-windup mechanism used in the implementation.

Writing the plant as \(G(z)=b(z)/a(z)\) and the controller as
\(G_C(z)=d(z)/c(z)\), the closed-loop transfer function is
\begin{equation}
  T(z)=\frac{b(z)d(z)}{q(z)} ,
  \label{eq:closed-loop-transfer}
\end{equation}
which is the form used throughout this report. The transient cannot be shaped
by the pole locations alone, because the numerator of
\eqref{eq:closed-loop-transfer} contains the controller zero
\begin{equation}
  z_0=-\frac{d_0}{d_1}.
  \label{eq:controller-zero}
\end{equation}
This is the same structural effect described in the reader as a phase error
caused by the different numerator order. A one-degree-of-freedom controller
provides no independent numerator design and cannot remove the zero without
changing the controller structure.
